\documentclass[a4paper,11pt]{article}

% ─── Geometry ────────────────────────────────────────────────────────────────
\usepackage{geometry}
\geometry{margin=1in, headheight=15pt}

% ─── Mathematics ─────────────────────────────────────────────────────────────
\usepackage{amsmath}
\usepackage{amssymb}
\usepackage{mathtools}

% ─── Graphics / Colours ──────────────────────────────────────────────────────
\usepackage{graphicx}
\usepackage{xcolor}
\definecolor{blue3}{RGB}{0,0,180}
\definecolor{green3}{RGB}{0,120,0}
\definecolor{orange3}{RGB}{200,100,0}
\definecolor{red3}{RGB}{180,0,0}
\definecolor{grey3}{RGB}{80,80,80}

% ─── Units ───────────────────────────────────────────────────────────────────
\usepackage{siunitx}
\sisetup{detect-all}

% ─── Links (hyperref BEFORE cleveref) ────────────────────────────────────────
\usepackage[colorlinks=true,linkcolor=blue3,citecolor=blue3,urlcolor=blue3]{hyperref}
\usepackage{cleveref}

% ─── Boxes ───────────────────────────────────────────────────────────────────
\usepackage{tcolorbox}
\tcbuselibrary{skins,breakable}

\newtcolorbox{pointsbox}[1][]{
    colback=grey3!8,
    colframe=grey3,
    fonttitle=\bfseries\small,
    #1
}

\newtcolorbox{hintbox}[1][]{
    colback=orange3!8,
    colframe=orange3,
    title={Hint:},
    fonttitle=\bfseries,
    #1
}

% ─── Tables / Lists ──────────────────────────────────────────────────────────
\usepackage{booktabs}
\usepackage{enumitem}

% ─── Notation commands ───────────────────────────────────────────────────────
\newcommand{\ktilde}{\tilde{k}}
\newcommand{\ytilde}{\tilde{y}}
\newcommand{\ctilde}{\tilde{c}}
\newcommand{\dd}{\mathrm{d}}
\newcommand{\bgp}{\mathrm{BGP}}

% ─── Header / Footer ─────────────────────────────────────────────────────────
\usepackage{fancyhdr}
\pagestyle{fancy}
\fancyhf{}
\lhead{\textbf{14E022 Macroeconomics I} \quad Practice Exam 3}
\rhead{BSE 2025--26}
\cfoot{\thepage}

% ═════════════════════════════════════════════════════════════════════════════
\begin{document}
% ═════════════════════════════════════════════════════════════════════════════

\begin{center}
    {\LARGE \textbf{14E022 Macroeconomics I}}\\[6pt]
    {\Large \textbf{Practice Exam 3}}\\[8pt]
    {\large BSE --- Academic Year 2025--26}
\end{center}

\vspace{8pt}
\noindent\textit{Instructions: the exam is made of 3 problems and a set of short questions totalling \textbf{100 points}.
You have \textbf{2 hours}. Show all derivations clearly. For short questions, concise and precise answers score highest.}
\vspace{4pt}

\noindent\rule{\textwidth}{0.4pt}

% ─────────────────────────────────────────────────────────────────────────────
\section{Question 1 --- Romer (1990) Expanding Varieties \hfill (25 points)}
% ─────────────────────────────────────────────────────────────────────────────

Consider a three-sector economy consisting of a final goods sector, an intermediate goods sector,
and an R\&D sector.  The representative final goods firm produces output using production workers
$L_Y$ and a continuum of differentiated intermediate inputs:
\[
    Y_t = L_Y^{1-\alpha} \int_0^{M_t} x_i^{\alpha}\,\dd i,
    \qquad \alpha \in (0,1),
\]
where $M_t$ is the total mass of available varieties and $x_i$ is the quantity of variety $i$
used.  Each intermediate variety is produced one-for-one from the final good (marginal cost
$= 1$).  The R\&D technology expands the variety frontier at rate
\[
    \dot{M}_t = \lambda M_t L_{A,t},
\]
where $L_{A,t}$ is the mass of workers employed in R\&D and $\lambda > 0$ is a productivity
parameter.  Households have CRRA preferences with risk-aversion $\sigma > 0$, discount rate
$\rho > 0$, and there is no population growth.  Total labour supply is $L = L_Y + L_A$.

\begin{enumerate}[label=\textbf{\arabic*.}, leftmargin=*, itemsep=6pt]

    \item \textbf{(5 points)} Final goods sector FOC.

    \begin{enumerate}[label=(\alph*), itemsep=4pt]
        \item Set up the final goods firm's profit-maximisation problem taking output as
        the num\'{e}raire and each $p_i$ as given.  Derive the isoelastic inverse demand
        for variety $i$:
        \[
            p_i = \alpha L_Y^{1-\alpha}\,x_i^{\alpha-1}.
        \]
        \item Invert this to obtain the demand function $x_i(p_i)$ and show that the price
        elasticity of demand is $\varepsilon = 1/(1-\alpha) > 1$.
    \end{enumerate}

    \item \textbf{(6 points)} Intermediate monopolist's pricing and profits.

    Each variety $i$ is supplied by a monopolist who bought the blueprint from the R\&D sector.

    \begin{enumerate}[label=(\alph*), itemsep=4pt]
        \item Using the demand function from part 1 and the CES markup formula
        $p = MC \cdot \varepsilon/(\varepsilon - 1)$, derive the monopoly price
        $p_i^* = 1/\alpha$.  (You may alternatively derive it directly from the
        monopolist's FOC.)
        \item Substitute $p_i^* = 1/\alpha$ into the demand function to obtain the
        equilibrium quantity:
        \[
            x_i^* = \alpha^{2/(1-\alpha)}\,L_Y.
        \]
        \item Show that equilibrium flow profit is
        \[
            \pi = (p_i^* - 1)\,x_i^* = (1-\alpha)\,\alpha^{\alpha/(1-\alpha)}\,L_Y.
        \]
        Verify that $\pi = \tfrac{1-\alpha}{\alpha}\,x_i^*$.
    \end{enumerate}

    \item \textbf{(7 points)} R\&D equilibrium and the BGP growth rate.

    \begin{enumerate}[label=(\alph*), itemsep=4pt]
        \item \textbf{No-arbitrage (Bellman equation).}  Let $V_t$ be the value of owning
        one blueprint (one variety).  Unlike quality ladders, there is no creative
        destruction in Romer: once created, a variety exists forever.  Write the
        no-arbitrage condition $r V = \pi + \dot{V}$.  On the BGP where $V$ is constant,
        show this simplifies to $r V = \pi$.

        \item \textbf{R\&D free entry.}  A researcher can create a new variety by hiring
        labour at wage $w$.  The cost of discovering one new variety is $w/(\lambda M)$
        (the reciprocal of the marginal product of research labour).  Free entry drives
        profits from R\&D to zero:
        \[
            V = \frac{w}{\lambda M}.
        \]
        Combine with the no-arbitrage condition to show $r = \lambda M \pi / w$.

        \item \textbf{Labour market and Euler equation.}  The wage equals the marginal
        product of $L_Y$ in the final sector: $w = (1-\alpha)Y/L_Y$.  Use the resource
        allocations $L_Y + L_A = L$ and the Euler equation $\dot{\ctilde}/\ctilde =
        (r-\rho)/\sigma$ to derive the BGP growth rate:
        \[
            g^{\mathrm{mkt}} = \frac{\lambda(1-\alpha)L - \rho}{\sigma + 1}.
        \]
        Identify the \textbf{scale effect}: explain why $g^{\mathrm{mkt}}$ rises with total
        labour supply $L$.
    \end{enumerate}

    \begin{hintbox}
        On the BGP, $L_Y$ and $L_A$ are constant fractions of $L$.  Use the free-entry
        condition and wage expression to pin down $L_Y^*$, substitute into the Euler
        equation.
    \end{hintbox}

    \item \textbf{(7 points)} Social planner vs.\ market.

    \begin{enumerate}[label=(\alph*), itemsep=4pt]
        \item Explain the \textbf{two distortions} in the decentralised equilibrium:
        \begin{enumerate}[label=(\roman*), itemsep=2pt]
            \item \emph{Markup distortion}: the monopolist sets $p_i = 1/\alpha > 1 = MC$,
            so each variety is under-used relative to the social optimum.  How does the
            planner set quantity?
            \item \emph{Knowledge spillover}: when a new variety is created, the inventor
            captures only the discounted flow of monopoly profits $V$, but the social value
            also includes consumer surplus and the indirect effect of variety expansion on
            the productivity of future R\&D ($\partial \dot{M}/\partial M$).  State the
            direction of each distortion (under- or over-provision).
        \end{enumerate}
        \item Show that $g^{\mathrm{plnr}} > g^{\mathrm{mkt}}$.  You do not need to
        derive $g^{\mathrm{plnr}}$ explicitly; it suffices to argue that the planner uses
        more of each variety (correcting the markup) and internalises the spillover,
        both of which raise the incentive to conduct R\&D.

        \item Identify the \textbf{optimal policy mix} that implements the first-best:
        \begin{enumerate}[label=(\roman*), itemsep=2pt]
            \item A subsidy $s_x = 1-\alpha$ to intermediate goods purchasers (final goods
            firms pay only $\alpha$ per unit, eliminating the markup wedge).
            \item An R\&D subsidy that internalises the knowledge-spillover externality.
        \end{enumerate}
        Argue that a single instrument \emph{cannot} achieve the first-best because the
        two distortions are independent.
    \end{enumerate}

\end{enumerate}

\clearpage

% ─────────────────────────────────────────────────────────────────────────────
\section{Question 2 --- Structural Change: Baumol Effect and Non-Homothetic Preferences \hfill (25 points)}
% ─────────────────────────────────────────────────────────────────────────────

Consider a two-sector economy with manufacturing ($M$) and services ($S$).  Production
technologies are linear in labour:
\[
    Y_M = A_M L_M, \qquad Y_S = A_S L_S,
\]
with productivity dynamics $\dot{A}_M/A_M = g > 0$ and $\dot{A}_S/A_S = 0$ (services are
the stagnant sector).  Labour markets are competitive and the total labour force is fixed:
$L_M + L_S = L$.

\begin{enumerate}[label=\textbf{\arabic*.}, leftmargin=*, itemsep=6pt]

    \item \textbf{(5 points)} Competitive factor markets and relative prices.

    \begin{enumerate}[label=(\alph*), itemsep=4pt]
        \item With a competitive labour market, wages are equalised across sectors.  Show
        that profit maximisation in each sector implies $p_M A_M = p_S A_S = w$.  Hence
        derive:
        \[
            \frac{p_S}{p_M} = \frac{A_M}{A_S}.
        \]
        \item Explain why this relative price grows at rate $g$ over time.  What is the
        economic interpretation?  This is the \textbf{cost disease} (Baumol, 1967):
        wages in the stagnant sector must keep pace with the productive sector,
        pushing up services prices.
    \end{enumerate}

    \item \textbf{(7 points)} Baumol mechanism with CES preferences.

    Households have CES preferences:
    \[
        U = \bigl[\omega\,Y_M^{\rho} + (1-\omega)\,Y_S^{\rho}\bigr]^{1/\rho},
        \qquad \rho \in (-\infty, 1)\setminus\{0\},
    \]
    with elasticity of substitution $\sigma_s = 1/(1-\rho)$.

    \begin{enumerate}[label=(\alph*), itemsep=4pt]
        \item Derive the expenditure-optimal consumption ratio from the household's static
        optimum:
        \[
            \frac{Y_S}{Y_M} = \left[\frac{(1-\omega)}{\omega}\cdot\frac{p_M}{p_S}
            \right]^{\sigma_s}.
        \]
        \item Using $Y_j = A_j L_j$, convert the output ratio to an employment ratio:
        \[
            \frac{L_S}{L_M} = \left(\frac{1-\omega}{\omega}\right)^{\sigma_s}
            \left(\frac{A_M}{A_S}\right)^{1-\sigma_s}.
        \]
        \item Analyse the dynamics of $L_S/L_M$ over time (as $A_M$ grows):
        \begin{enumerate}[label=(\roman*), itemsep=2pt]
            \item If $\sigma_s = 1$ (Cobb-Douglas): show employment shares are
            \textbf{constant} despite rising $p_S$.
            \item If $\sigma_s < 1$ (complements): show $L_S/L_M$ rises over time,
            so labour flows to the stagnant sector (\textbf{Baumol result}).
            \item If $\sigma_s > 1$ (substitutes): show labour flows to the productive
            sector.
        \end{enumerate}
        \item Show that aggregate productivity growth satisfies $\gamma_Y = g\cdot(L_M/L)$
        and explain why structural change implies $\gamma_Y \to 0$ as $L_M/L \to 0$
        (the Baumol growth drag).
    \end{enumerate}

    \item \textbf{(6 points)} Non-homothetic preferences.

    Suppose instead that households have Stone-Geary utility:
    \[
        U = \log(Y_M - \bar{c}_M) + \log(Y_S), \qquad \bar{c}_M > 0,
    \]
    where $\bar{c}_M$ is a subsistence level of manufactured consumption.  Total
    nominal income is $I = p_M Y_M + p_S Y_S$.

    \begin{enumerate}[label=(\alph*), itemsep=4pt]
        \item Derive the demand functions by maximising $U$ subject to the budget
        constraint.  Show:
        \[
            Y_M = \bar{c}_M + \frac{I - p_M \bar{c}_M}{2\,p_M},
            \qquad
            Y_S = \frac{I - p_M \bar{c}_M}{2\,p_S}.
        \]
        \item Derive the \textbf{income elasticity of demand} for manufacturing:
        \[
            \varepsilon_M^I = \frac{\partial Y_M}{\partial I}\cdot\frac{I}{Y_M}
            = 1 - \frac{\bar{c}_M p_M}{I - \bar{c}_M p_M} \cdot \frac{p_M \bar{c}_M}{p_M Y_M} \,.
        \]
        Show that $\varepsilon_M^I < 1$, i.e.\ manufacturing is a \textbf{necessity}.
        \item As income $I$ rises (with prices fixed), show that the expenditure share of
        manufacturing falls.  This is \textbf{demand-side structural change} that occurs
        even with equal productivity growth across sectors.
    \end{enumerate}

    \item \textbf{(7 points)} Synthesis and empirical comparison.

    \begin{enumerate}[label=(\alph*), itemsep=4pt]
        \item In the Baumol model (part 2), structural change occurs even with
        \textbf{homothetic preferences} provided $\sigma_s < 1$.  What is the key
        supply-side mechanism?  Why is the elasticity of substitution the critical
        parameter?

        \item In the non-homothetic model (part 3), structural change occurs even with
        \textbf{equal productivity growth} across sectors.  What is the key demand-side
        mechanism?  What does the income elasticity below unity imply for how households
        allocate their marginal dollar as they grow richer?

        \item Duarte and Restuccia (2010) document average annual TFP growth rates of
        approximately 3.8\% in agriculture, 2.4\% in manufacturing, and 1.3\% in services
        over 1956--2004 across a broad panel of countries.  Is this differential consistent
        with the Baumol cost-disease mechanism?  In the data, employment shares in services
        rise and manufacturing employment falls as countries grow richer.  Based on this
        evidence, discuss the relative importance of supply-side (Baumol) versus
        demand-side (non-homothetic) channels for understanding structural change.
    \end{enumerate}

\end{enumerate}

\clearpage

% ─────────────────────────────────────────────────────────────────────────────
\section{Question 3 --- Solow Model: Golden Rule, Convergence Speed, and Dynamic Inefficiency \hfill (25 points)}
% ─────────────────────────────────────────────────────────────────────────────

Consider the standard Solow growth model in continuous time.  Aggregate production is
\[
    Y_t = K_t^{\alpha}(A_t L_t)^{1-\alpha}, \qquad \alpha \in (0,1),
\]
with saving rate $s \in (0,1)$, population growth $\dot{L}/L = n > 0$, TFP growth
$\dot{A}/A = g > 0$, and capital depreciation $\delta > 0$.

\begin{enumerate}[label=\textbf{\arabic*.}, leftmargin=*, itemsep=6pt]

    \item \textbf{(5 points)} Per-effective-worker dynamics.

    \begin{enumerate}[label=(\alph*), itemsep=4pt]
        \item Define per-effective-worker variables $\ktilde \equiv K/(AL)$ and
        $\ytilde \equiv Y/(AL) = \ktilde^{\alpha}$.  Derive the law of motion:
        \[
            \dot{\ktilde} = s\,\ktilde^{\alpha} - (n+g+\delta)\,\ktilde.
        \]
        \item Show using the Inada conditions ($f'(\ktilde) \to \infty$ as $\ktilde \to 0$,
        $f'(\ktilde) \to 0$ as $\ktilde \to \infty$) that a unique interior steady state
        $\ktilde^* > 0$ exists.  Solve for it:
        \[
            \ktilde^* = \left[\frac{s}{n+g+\delta}\right]^{\!1/(1-\alpha)}.
        \]
    \end{enumerate}

    \item \textbf{(7 points)} Golden Rule capital accumulation.

    \begin{enumerate}[label=(\alph*), itemsep=4pt]
        \item At the steady state $\ktilde^*$, express per-effective-worker consumption as:
        \[
            \ctilde^* = f(\ktilde^*) - (n+g+\delta)\,\ktilde^* = \ktilde^{*\alpha}
            - (n+g+\delta)\,\ktilde^*.
        \]
        (Note: $\ctilde^* = (1-s)\ktilde^{*\alpha}$; both expressions are equivalent at
        the steady state.)

        \item To find the \textbf{Golden Rule saving rate} that maximises $\ctilde^*$,
        treat $\ktilde^*$ as a free variable (each $s$ pins down a different $\ktilde^*$).
        Differentiate $\ctilde^*$ with respect to $\ktilde^*$ and set to zero:
        \[
            \frac{\partial \ctilde^*}{\partial \ktilde^*} = 0
            \;\Longrightarrow\;
            f'(\ktilde^*_{\mathrm{GR}}) = n+g+\delta.
        \]
        For Cobb-Douglas, solve explicitly: $\alpha\,(\ktilde^*_{\mathrm{GR}})^{\alpha-1}
        = n+g+\delta$.  Show that the corresponding saving rate is $s_{\mathrm{GR}} = \alpha$.

        \item Interpret the Golden Rule.  What happens to steady-state consumption if
        $s > s_{\mathrm{GR}}$?  If $s < s_{\mathrm{GR}}$?  Define the concept of
        \textbf{dynamic inefficiency} (over-saving) and explain why a Pareto improvement
        is possible when $s > s_{\mathrm{GR}}$.
    \end{enumerate}

    \item \textbf{(6 points)} Speed of convergence.

    \begin{enumerate}[label=(\alph*), itemsep=4pt]
        \item Linearise the law of motion $\dot{\ktilde} = s\ktilde^{\alpha}
        - (n+g+\delta)\ktilde$ around $\ktilde^*$.  Define the log-deviation
        $\hat{k}_t \equiv \ln(\ktilde_t/\ktilde^*)$ and show that:
        \[
            \dot{\hat{k}}_t \approx -\beta\,\hat{k}_t,
            \qquad \text{where} \quad
            \beta = (1-\alpha)(n+g+\delta).
        \]
        This implies $\hat{k}_t = \hat{k}_0\,e^{-\beta t}$: the gap to steady state
        shrinks exponentially at rate $\beta$.

        \item Calibrate the formula.  With $\alpha = 1/3$, $n = 0.01$, $g = 0.02$,
        $\delta = 0.05$, compute $\beta$.  The empirical estimate from cross-country
        regressions (Mankiw, Romer, Weil 1992) is approximately 2\% per year.
        Why is the empirical estimate slower than the basic calibration?
        (\textit{Hint}: human capital augmentation effectively raises the capital share
        toward $2/3$.)

        \item How does $\beta$ depend on $\alpha$?  Provide economic intuition: why does
        a higher capital share imply \emph{slower} convergence to the steady state?
    \end{enumerate}

    \item \textbf{(7 points)} Dynamic inefficiency.

    \begin{enumerate}[label=(\alph*), itemsep=4pt]
        \item Define \textbf{dynamic inefficiency} formally: the steady state is
        dynamically inefficient if $r^* < n + g$, where $r^* = \alpha\ktilde^{*\alpha-1}
        - \delta$ is the net return to capital.  Show that this condition is equivalent
        to $s > s_{\mathrm{GR}} = \alpha$ using the expression for $\ktilde^*$ and the
        fact that $r^* K = \alpha Y$ at the steady state.

        \item \textbf{Abel et al.\ (1989) test.}  A sufficient condition for dynamic
        \emph{efficiency} in a stochastic economy is that gross capital income exceeds
        gross investment:
        \[
            r^* K \;\geq\; \delta K + \dot{K},
            \quad\text{i.e.,}\quad
            \alpha Y \;\geq\; (n+g+\delta)K.
        \]
        Abel et al.\ document that this inequality holds for all G7 economies throughout
        the post-war period.  What does this imply about the empirical relevance of
        dynamic inefficiency?

        \item Can the \textbf{Ramsey model} be dynamically inefficient?  Explain
        rigorously.  Specifically, show that the Transversality Condition (TVC):
        \[
            \lim_{t\to\infty} e^{-\rho t}\,\lambda_t\,k_t = 0
        \]
        implies $r^* > n + g$ on the BGP (where $\lambda_t$ is the costate on capital
        in per-capita terms and $\rho > n$ by assumption for bounded utility).  Contrast
        this with the Solow model, where $s$ is exogenous and may inadvertently push the
        economy into the over-saving region.
    \end{enumerate}

\end{enumerate}

\clearpage

% ─────────────────────────────────────────────────────────────────────────────
\section{Short Questions \hfill (25 points)}
% ─────────────────────────────────────────────────────────────────────────────

\noindent\textit{For these questions, short and sharp answers will receive the highest marks.
Aim for 3--6 sentences per sub-question unless otherwise indicated.}

\vspace{8pt}

\begin{enumerate}[label=\textbf{\arabic*.}, leftmargin=*, itemsep=14pt]

    % ── SQ1 ──────────────────────────────────────────────────────────────────
    \item \textbf{(9 points) Kaldor Facts and Balanced Growth.}

    \begin{enumerate}[label=(\alph*), itemsep=6pt]
        \item State \textbf{five} Kaldor (1961) stylized facts about long-run economic
        growth in developed economies.  Which family of production functions is consistent
        with all five facts on a Balanced Growth Path (BGP)?  Explain \emph{why} this
        family is special.

        \item Show formally that a Cobb-Douglas production function
        $Y = K^{\alpha}(AL)^{1-\alpha}$ implies a \textbf{constant capital income share}
        $\alpha = rK/Y$ regardless of the levels of $K$, $L$, or $A$.  Relate this
        property to why Solow + Cobb-Douglas generates a BGP with constant growth rates
        of all per-worker variables.

        \item Karabarbounis and Neiman (2014) document a declining global labour share
        since the 1980s.  Does this violation of the constant-shares Kaldor fact
        invalidate the Solow framework?  Briefly describe one model modification
        (e.g.\ a non-unit elasticity of substitution between capital and labour) that
        could accommodate a trending labour share while preserving balanced growth.
    \end{enumerate}

    % ── SQ2 ──────────────────────────────────────────────────────────────────
    \item \textbf{(8 points) Knowledge Spillovers and the $\phi$ Parameter.}

    Consider the general R\&D technology
    \[
        \dot{M}_t = \lambda\,M_t^{\phi}\,L_{A,t}, \qquad \lambda > 0.
    \]

    \begin{enumerate}[label=(\alph*), itemsep=6pt]
        \item Derive the BGP growth rate $g_M \equiv \dot{M}/M$ for three cases:
        (i)~$\phi = 1$ (Romer 1990); (ii)~$\phi < 1$ (Jones 1995 semi-endogenous growth);
        (iii)~$\phi = 0$ (no knowledge spillovers).
        In each case, what is the role of \textbf{population growth} $n$ in sustaining
        positive long-run growth in $M$?

        \item ``\textit{Standing on shoulders}'' (past knowledge raises the productivity
        of current researchers) vs.\ ``\textit{fishing out}'' (the stock of easy
        discoveries is depleted by past R\&D):  which value of $\phi$ corresponds to
        each metaphor?  What empirical evidence from patent citation studies informs
        the plausible range of $\phi$?

        \item An R\&D subsidy raises $L_A$.  In which case(s) does this permanently
        raise the \textbf{long-run growth rate} $g_M$?  In which case(s) does it affect
        only the \textbf{level} of the BGP path (transitional effect)?  State the key
        equation that distinguishes the two outcomes.
    \end{enumerate}

    % ── SQ3 ──────────────────────────────────────────────────────────────────
    \item \textbf{(8 points) Consumption Tax vs.\ Capital Income Tax in the Ramsey Model.}

    Consider a Ramsey economy where the government levies a \textbf{proportional
    consumption tax} $\tau_c \in (0,1)$, returning all revenue as a lump-sum transfer $T_t$
    to households.  The household's flow budget constraint is:
    \[
        \dot{a}_t = r_t a_t + w_t + T_t - (1+\tau_c)\,c_t,
    \]
    where $a_t$ denotes assets, and households maximise $\int_0^{\infty} e^{-\rho t}
    u(c_t)\,\dd t$ with $u(c) = c^{1-\sigma}/(1-\sigma)$.

    \begin{enumerate}[label=(\alph*), itemsep=6pt]
        \item Set up the \emph{current-value} Hamiltonian with costate variable $\mu_t$.
        Derive the optimality conditions and eliminate $\mu_t$ to obtain the
        \textbf{Euler equation} for $\dot{c}/c$.

        \item Show that the consumption tax \textbf{does not distort} the intertemporal
        allocation of consumption: the Euler equation is \emph{identical} to the no-tax
        case.  Provide intuition: why does taxing consumption at a constant proportional
        rate leave the intertemporal price ratio $c_{t+\dd t}/c_t$ unchanged?

        \item Contrast with the \textbf{capital income tax} $\tau_K$ studied in
        Practice Exam 1, which replaces $r$ with $r(1-\tau_K)$ in the Euler equation.
        Explain the asymmetry: why does $\tau_K$ distort saving while $\tau_c$ does not,
        even though both taxes reduce household resources?  What is the implication of
        this asymmetry for \textbf{optimal tax theory} (e.g., the Chamley-Judd result)?
    \end{enumerate}

\end{enumerate}

\vspace{20pt}
\noindent\rule{\textwidth}{0.4pt}

\begin{center}
    \textit{End of Practice Exam 3. Good luck!}
\end{center}

% ─────────────────────────────────────────────────────────────────────────────
\clearpage
\appendix
\section*{Appendix: Key Formulas Permitted in Exam}
% ─────────────────────────────────────────────────────────────────────────────

The following results may be used without re-derivation unless the question explicitly asks
you to derive them.

\subsection*{A.1 ~ Log-differentiation}
If $Z_t = X_t^a Y_t^b$, then $\gamma_Z = a\gamma_X + b\gamma_Y$.

\subsection*{A.2 ~ Ramsey / Euler equation (standard, no tax)}
\[
    \frac{\dot{c}}{c} = \frac{1}{\sigma}(r - \rho),
    \qquad
    \frac{\dot{\ctilde}}{\ctilde} = \frac{1}{\sigma}(r - \delta - \rho - \sigma x),
\]
where $r = f'(\ktilde) - \delta$, $x = \dot{A}/A$, and $n$ is population growth.

\subsection*{A.3 ~ Hamiltonian first-order conditions (current-value)}
Control $u$, state $x$, costate $\lambda$, discount $\rho$:
\[
    H_u = 0, \qquad \dot{\lambda} = \rho\lambda - H_x.
\]

\subsection*{A.4 ~ Romer (1990) demand function and monopoly markup}
Final sector: $Y = L_Y^{1-\alpha}\int_0^M x_i^\alpha\,\dd i$.
\[
    \text{Demand for variety } i{:}\quad x_i = \left(\frac{\alpha L_Y^{1-\alpha}}{p_i}\right)^{1/(1-\alpha)}, \qquad
    \varepsilon = \frac{1}{1-\alpha}.
\]
\[
    p_i^* = \frac{1}{\alpha}, \qquad
    x_i^* = \alpha^{2/(1-\alpha)}L_Y, \qquad
    \pi = (1-\alpha)\alpha^{\alpha/(1-\alpha)}L_Y.
\]

\subsection*{A.5 ~ Romer BGP growth rate}
\[
    g^{\mathrm{mkt}} = \frac{\lambda(1-\alpha)L - \rho}{\sigma + 1}.
\]
Scale effect: $\partial g^{\mathrm{mkt}}/\partial L > 0$.

\subsection*{A.6 ~ Solow steady state and Golden Rule}
\[
    \ktilde^* = \left[\frac{s}{n+g+\delta}\right]^{1/(1-\alpha)}, \qquad
    s_{\mathrm{GR}} = \alpha \quad (\text{Cobb-Douglas}).
\]

\subsection*{A.7 ~ Speed of convergence}
\[
    \beta = (1-\alpha)(n+g+\delta).
\]
Interpretation: the gap $\hat{k}_t = \ln(\ktilde_t/\ktilde^*)$ closes at rate $\beta$ per year.

\subsection*{A.8 ~ Abel et al.\ (1989) dynamic efficiency test}
Economy is dynamically efficient if gross capital income $\geq$ gross investment:
\[
    \alpha Y \;\geq\; (n+g+\delta)K.
\]

\subsection*{A.9 ~ Stone-Geary demand functions}
Utility $U = \log(Y_M - \bar{c}_M) + \log Y_S$, income $I = p_M Y_M + p_S Y_S$:
\[
    Y_M = \bar{c}_M + \frac{I - p_M\bar{c}_M}{2p_M},
    \qquad
    Y_S = \frac{I - p_M\bar{c}_M}{2p_S}.
\]

\subsection*{A.10 ~ General R\&D technology ($\phi$ parameter)}
$\dot{M} = \lambda M^{\phi} L_A$.  On the BGP:
\[
    \phi = 1{:}\quad g_M = \lambda L_A; \qquad
    \phi < 1{:}\quad g_M = \frac{n}{1-\phi}; \qquad
    \phi = 0{:}\quad g_M = \lambda L_A \text{ (requires growing $L_A$)}.
\]

\end{document}
